\begingroup
\renewcommand{\clearpageforsection}{}
\exercisesheader{}
\endgroup

% 15

\eoce{\qt{Mengidentifikasi hipotesis, Bagian I\label{
}}
Tuliskan hipotesis nol dan hipotesis alternatif dengan kata-kata, lalu dengan simbol,
untuk masing-masing situasi berikut.
\begin{parts}
\item
    Sebuah perusahaan bimbingan belajar ingin mengetahui apakah sebagian besar
    siswa cenderung mengalami peningkatan nilai (atau tidak) setelah
    menggunakan layanannya.
    Perusahaan tersebut mengambil sampel sebanyak 200 siswa yang menggunakan layanannya
    selama setahun terakhir dan menanyakan apakah nilai mereka
    meningkat atau menurun dibandingkan dengan tahun sebelumnya.
\item
    Pimpinan sebuah perusahaan mengkhawatirkan dampak March Madness,
    kejuaraan bola basket yang diadakan setiap musim semi di AS, terhadap produktivitas
    karyawan.
    Mereka mengestimasi bahwa pada hari kerja biasa, karyawan menggunakan rata-rata
    15 menit waktu kerja untuk memeriksa surel pribadi,
    melakukan panggilan telepon pribadi, dan sebagainya.
    Mereka juga mengumpulkan data tentang banyaknya waktu kerja yang digunakan karyawan
    untuk kegiatan nonpekerjaan tersebut selama March Madness.
    Mereka ingin menentukan apakah data ini memberikan bukti yang meyakinkan
    bahwa produktivitas karyawan berubah selama March Madness.
\end{parts}
}{}

% 16

\eoce{\qt{Mengidentifikasi hipotesis, Bagian II\label{identify_hypotheses_prop_and_mean_2}} 
Tuliskan hipotesis nol dan hipotesis alternatif dengan kata-kata dan menggunakan simbol
untuk masing-masing situasi berikut.
\begin{parts}
\item
    Sejak 2008, jaringan restoran di California diwajibkan
    mencantumkan jumlah kalori setiap hidangan pada menu. Sebelum menu
    mencantumkan jumlah kalori, rata-rata asupan kalori pengunjung
    di sebuah restoran adalah 1100 kalori.
    Setelah jumlah kalori mulai dicantumkan pada menu,
    seorang ahli gizi mengumpulkan data tentang jumlah kalori yang dikonsumsi
    di restoran tersebut dari sampel acak para pengunjung.
    Apakah data ini memberikan bukti yang meyakinkan bahwa terdapat perbedaan pada
    rata-rata asupan kalori para pengunjung di restoran tersebut?
\item
    Negara bagian Wisconsin ingin mengetahui
    proporsi penduduk dewasanya yang mengonsumsi alkohol
    selama setahun terakhir,
    khususnya apakah proporsi tersebut berbeda dari
    proporsi nasional sebesar 70\%.
    Untuk membantu menjawab pertanyaan ini, mereka mengambil
    sampel acak sebanyak 852 penduduk dan menanyakan
    konsumsi alkohol mereka.
\end{parts}
}{}

% 17

\eoce{\qt{Komunikasi daring\label{online_communication_prop_ht_errors}}
Sebuah studi mengemukakan bahwa 60\% mahasiswa menghabiskan
10~jam atau lebih per minggu untuk berkomunikasi dengan orang lain secara daring.
Anda yakin bahwa pernyataan ini tidak benar dan memutuskan untuk mengambil
sampel sendiri guna melakukan uji hipotesis.
Anda mengambil sampel acak sebanyak 160 mahasiswa dari asrama Anda
dan mendapati bahwa 70\% di antaranya menghabiskan 10~jam atau lebih per minggu
untuk berkomunikasi dengan orang lain secara daring.
Seorang~teman Anda, yang menawarkan bantuan untuk melakukan
uji hipotesis tersebut, menyusun pasangan
hipotesis berikut.
Tunjukkan setiap kesalahan yang Anda temukan.
\begin{align*}
H_0&: \hat{p} < 0.6 \\
H_A&: \hat{p} > 0.7
\end{align*}
}{}

% 18

\eoce{\qt{Menikah pada usia 25 tahun\label{married_at_25_prop_ht_errors}}
Sebuah studi mengemukakan bahwa 25\% orang berusia 25 tahun telah
menikah.
Anda yakin bahwa pernyataan ini tidak benar dan memutuskan untuk mengambil
sampel sendiri guna melakukan uji hipotesis.
Dari data sensus, Anda mengambil sampel acak
orang berusia 25 tahun sebanyak 776 orang dan mendapati bahwa
24\% di antaranya telah menikah.
Seorang teman Anda menawarkan bantuan untuk menyusun
uji hipotesis dan menghasilkan pasangan
hipotesis berikut.
Tunjukkan setiap kesalahan yang Anda temukan.
\begin{align*}
H_0&: \hat{p} = 0.24 \\
H_A&: \hat{p} \neq 0.24
\end{align*}
}{}

% 19

\eoce{\qt{Tingkat perundungan siber\label{cyberbullying_prop_ci_ht}}
Para remaja disurvei mengenai perundungan siber, dan
54\% hingga 64\% melaporkan pernah mengalami perundungan siber
(interval kepercayaan 95\%).\footfullcite{pew_cyber_bully_2018}
Jawablah pertanyaan-pertanyaan berikut berdasarkan interval tersebut.
\begin{parts}
\item 
    Sebuah surat kabar menyatakan bahwa mayoritas remaja
    pernah mengalami perundungan siber.
    Apakah klaim ini didukung oleh interval kepercayaan tersebut?
    Jelaskan alasan Anda.
\item\label{cyberbullying_prop_ci_ht_researcher}
    Seorang peneliti menduga bahwa 70\% remaja pernah
    mengalami perundungan siber.
    Apakah dugaan ini didukung oleh interval kepercayaan tersebut?
    Jelaskan alasan Anda.
\item
    Tanpa benar-benar menghitung intervalnya, tentukan
    apakah dugaan peneliti pada
    bagian~(\ref{cyberbullying_prop_ci_ht_researcher})
    akan didukung berdasarkan interval kepercayaan 90\%.
\end{parts}
}{}

% Boundary-clean reader removes a legacy forced page break.

% 20

\eoce{\qt{Menunggu di IGD, Bagian II\label{er_wait_ci_ht_prop_ok}}
Latihan~\ref{er_wait_intro_prop_ok} 
memberikan interval kepercayaan 95\% untuk rata-rata waktu tunggu
di instalasi gawat darurat (IGD), yaitu (128 menit, 147 menit).
Jawablah pertanyaan-pertanyaan berikut berdasarkan interval tersebut.
\begin{parts}
\item
    Sebuah surat kabar lokal menyatakan bahwa rata-rata waktu tunggu
    di IGD ini melebihi 3 jam.
    Apakah klaim ini didukung oleh interval kepercayaan tersebut?
    Jelaskan alasan Anda.
\item\label{er_wait_ci_ht_prop_ok_dean}
    Dekan Fakultas Kedokteran di rumah sakit ini menyatakan bahwa
    rata-rata waktu tunggunya adalah 2.2 jam.
    Apakah klaim ini didukung oleh interval kepercayaan tersebut?
    Jelaskan alasan Anda.
\item
    Tanpa benar-benar menghitung intervalnya,
    tentukan apakah klaim dekan pada
    bagian~(\ref{er_wait_ci_ht_prop_ok_dean})
    akan didukung berdasarkan interval kepercayaan 99\%.
\end{parts}
}{}

% 21

\eoce{\qt{Upah minimum, Bagian I\label{minimum_wage_prop_1}}
Apakah mayoritas orang dewasa AS meyakini bahwa kenaikan
upah minimum akan membantu perekonomian,
atau justru mayoritas tidak meyakininya?
Sebuah survei Rasmussen Reports terhadap sampel acak 1,000 orang dewasa AS mendapati
bahwa 42\% meyakini kenaikan itu akan membantu
perekonomian.\footfullcite{webpage:rasmussen-2019-raise-minimum-wage}
Lakukan uji hipotesis yang sesuai untuk membantu
menjawab pertanyaan penelitian tersebut.
}{}

% 22

\eoce{\qt{Tidur yang cukup\label{univ_students_enough_sleep}}
Sebanyak 400 mahasiswa diambil secara acak dari sebuah universitas besar,
dan 289 di antaranya mengatakan bahwa mereka tidak cukup tidur.
Lakukan uji hipotesis untuk memeriksa apakah hasil ini
merepresentasikan perbedaan yang signifikan secara statistik
dari 50\%, dengan menggunakan tingkat signifikansi 0.01.
}{}

% 23

\eoce{\qt{Menalar mundur, Bagian I\label{backwards_prop_1}}
Diberikan hipotesis-hipotesis berikut:
\begin{align*}
H_0&: p = 0.3 \\
H_A&: p \ne 0.3
\end{align*}
Diketahui ukuran sampelnya adalah 90.
Untuk proporsi sampel berapakah nilai-p sama dengan 0.05?
Asumsikan bahwa semua syarat yang diperlukan untuk inferensi telah terpenuhi.
}{}

% 24

\eoce{\qt{Menalar mundur, Bagian II\label{backwards_prop_2}}
Diberikan hipotesis-hipotesis berikut:
\begin{align*}
H_0&: p = 0.9 \\
H_A&: p \ne 0.9
\end{align*}
Diketahui ukuran sampelnya adalah 1,429.
Untuk proporsi sampel berapakah nilai-p sama dengan 0.01?
Asumsikan bahwa semua syarat yang diperlukan untuk inferensi telah terpenuhi.
}{}

% 25

\eoce{\qt{Menguji pengobatan fibromialgia\label{errors_fibromyalgia}} Seorang pasien bernama Diana
didiagnosis mengidap fibromialgia, sindrom nyeri tubuh jangka panjang, dan
diberi resep antidepresan. Karena bersikap skeptis, pada awalnya Diana tidak
percaya bahwa antidepresan akan membantu meredakan gejalanya. Namun, setelah
beberapa bulan menggunakan obat tersebut, ia menyimpulkan bahwa
antidepresan itu bekerja karena ia merasa gejalanya memang
semakin membaik.
\begin{parts}
\item Tuliskan hipotesis dengan kata-kata untuk pandangan skeptis Diana ketika ia
mulai menggunakan antidepresan.
\item Apa yang dimaksud dengan Galat Tipe~1 dalam konteks ini?
\item Apa yang dimaksud dengan Galat Tipe~2 dalam konteks ini?
\end{parts}
}{}

% 26

\eoce{\qtq{Manakah yang lebih besar\label{prop_which_higher_found_inf}}
Pada setiap bagian berikut terdapat suatu besaran yang menjadi perhatian dan dua
skenario (I dan II).
Untuk setiap bagian, nyatakan apakah besaran yang menjadi perhatian lebih besar
dalam skenario I, lebih besar dalam skenario II, atau sama besar
dalam kedua skenario.
\begin{parts}
\item
     Galat baku $\hat{p}$ ketika
     (I)~$n = 125$ atau (II)~$n = 500$.
\item
    Batas galat suatu interval kepercayaan
    ketika tingkat kepercayaannya adalah
    (I)~90\% atau (II)~80\%.
\item
    Nilai-p untuk statistik Z sebesar 2.5 yang dihitung
    berdasarkan (I)~sampel dengan $n = 500$ atau berdasarkan
    (II)~sampel dengan $n = 1000$.
\item
    Probabilitas terjadinya Galat Tipe~2 ketika
    hipotesis alternatif benar dan tingkat signifikansinya
    (I)~0.05 atau (II)~0.10.
\end{parts}
}{}
